\medskip\noindent
\textbf{A source decomposition that treats every nine-cell at once.}
We now use the target conclusions to extend the ordinary map,
keeping the choices made through dimension eight explicit.
Put $Q_i=K_i/K_4$ as in D.2.18--D.2.19, and work with the
two-local quotient cell models there. The five eight-cell
attachments in \textup{(D.2.19w)} give, by cell slides,
\[
 Q_8\simeq D'\vee M(\ZZ/2,7)_F\vee M(\ZZ/4,7)_J
                  \vee C_1^8\vee S^8_s,
 \tag{D.2.20s1}
\]
where
\[
\begin{split}
 D&=\operatorname{cofib}\left(
  \begin{pmatrix}4&\eta_5\\0&2\end{pmatrix}:
  S^5\vee S^6\longrightarrow S^5\vee S^6\right),\\
 D'&=D\cup_{\tau_d+\tau_e}e^8_H,\qquad
 C_1^8=\operatorname{cofib}\bigl((2,\eta_6):S^6\vee S^7\to S^6\bigr).
\end{split}
 \tag{D.2.20s2}
\]
Indeed slide the cell $G_8$ over $H_8$, leaving the
degree-two attachment on $F$; slide $B_0^2$ over $A_0C$,
leaving the degree-four attachment on $J$; and slide
$Q^5A_0$ over $A_0C$, leaving the free eight-sphere.
The integral orientation of the last sphere is
$s=Q^5A_0-A_0C$. The cell $A_0C$ supplies the top
cell of $C_1^8$, while $H_8$ supplies the top cell of $D'$.
These are actual comparisons for the mapping problem:
the quotients are four-connected, their dimensions are
at most eight, and Freudenthal realizes the stable cell
slides unstably. Maps into the two-local space $T$ may
then be pulled back to the original finite skeleta.

In D.2.19 the correction is zero on the source six-sphere
$z$ and on the free seven-spheres $F,J$. Its nullhomotopies
on the eight-cells may therefore be chosen so that the
correction is constant on all the factors of
\textup{(D.2.20s1)} except $D'$. This makes its restrictions
to the two Moore factors, $C_1^8$, and the last sphere
literally zero. On $D'$ it retains the Moore extension
$\mu$, the bottom class $\overline\alpha$, and the
six-sphere value $\rho$. Its extension across the
$H$-cell can be any of the allowed nullhomotopies;
we will not assume that its effect on the next
homotopy group is an arbitrary free choice.

\medskip\noindent
\textbf{The complete next group of the remaining factor.}
The two-local stable cofiber sequence defining $D$ gives
\[
 0\longrightarrow
   \ZZ/4\{v=i\nu_5\}\oplus\ZZ/2\{\tau_e\eta_7\}
 \longrightarrow\pi_8D\longrightarrow\ZZ/2\longrightarrow0.
 \tag{D.2.20s3}
\]
Here $\pi_3^S{}_{(2)}=\ZZ/8\{\nu\}$ and
$\eta^3=4\nu$; these are
\cite[\S2, (2.1)(c)--(d), p.~149]{I_Oka1986}.
On the preceding third-stem groups the defining matrix
has image generated by $4\nu$ in the bottom summand.
This gives the $\ZZ/4$ term and the surviving
second-stem class on the six-sphere. On the following
second-stem groups, the kernel is the bottom $\eta^2$
class. It is reached by $\tau_d\eta_7$: its top
pinch is $\eta_6\eta_7\ne0$. Since $\tau_d$ has
order two by \textup{(D.2.19p)}, this is an order-two
lift, splitting \textup{(D.2.20s3)}. Therefore
\[
 \pi_8D=
 \ZZ/4\{v\}\oplus
 \ZZ/2\{\tau_d\eta_7\}\oplus
 \ZZ/2\{\tau_e\eta_7\}.
 \tag{D.2.20s4}
\]

The attachment of $H$ kills the sum of the last two
classes. Its kernel on the preceding integral group
$\pi_7S^7$ consists of twice the top generator,
since $\tau_d+\tau_e$ has order two.
The resulting extension by this free rank-one kernel
splits as a group. We obtain the complete group
\[
 \pi_8D'=
 \ZZ_{(2)}\{\theta\}\oplus\ZZ/4\{v\}
                      \oplus\ZZ/2\{w\},
 \qquad
 h_{\ZZ_{(2)}}(\theta)=2H,\quad
 w=\tau_e\eta_7=\tau_d\eta_7.
 \tag{D.2.20s5}
\]
No canonical choice of $\theta$ is asserted: torsion
may be added to a lift of twice the $H$-cell. Crucially,
the class $v$ has order four. A primary fourth-square
detector would see only its parity and would leave
$2v$ undetected. Our argument retains both $2v$ and
$w$ in the group and checks their actual images.

All groups in \textup{(D.2.20s3)--(D.2.20s5)} are
also the unstable groups: $D,D'$ are four-connected
and the dimension used is eight. In particular,
$\theta$ is represented by an actual pointed sphere
map, not merely a stable class.

\medskip\noindent
\textbf{The correction kills the entire attaching group.}
Let $f:D'\to T$ be the chosen correction. Its two
torsion values are
\[
 f_*v=\overline\alpha\nu_5=0,\qquad
 f_*w=\rho\eta_6\eta_7=0.
 \tag{D.2.20s6}
\]
The first is a target conclusion of this theorem;
the second already follows from D.2.19. Thus the
previously undetected even multiple $2v$ is killed
as well.

The remaining value $\Theta=f_*\theta$ has zero
mod-two Hurewicz image. Indeed
$h_{\ZZ_{(2)}}(\theta)=2H$ implies
$h_{\FF_2}(\theta)=0$, independently of the choice
of its torsion lift or of the $H$-cell nullhomotopy.
The map $f$ lifts to $F=\fib(T\to B)$:
$\Map_*(D',K(\FF_2,2))$ is weakly contractible,
because $D'$ is four-connected. The inclusion
$H_8F\to H_8T$ is injective by the space splitting
$T\simeq B\times F$ of D.2.18. Naturality consequently
gives
\[
 h_{\FF_2}(\Theta)=0\quad\hbox{in }H_8F.
 \tag{D.2.20s7}
\]
The target result proved above says that the actual
mod-two Hurewicz homomorphism on $\pi_8F$ is
injective. It follows that $\Theta=0$.
This applies to every allowed filling of the $H$-cell;
no division in a stable group or unproved freedom
to choose $\Theta$ was used.

Equations \textup{(D.2.20s5)--(D.2.20s7)} show
that $f_*:\pi_8D'\to\pi_8T$ is zero. The
correction is zero on the other wedge factors in
\textup{(D.2.20s1)}. There are no cross Whitehead
terms in dimension eight, by the displayed
connectivities. Thus the correction kills all of
$\pi_8Q_8$. Every actual nine-cell attaching map,
including any whose even-$\nu$ component is not
determined by primary operations, has zero image.
Choose nullhomotopies on these cells to extend
the correction to $Q_9$.

Finally, the canonical summand already constructed
in \textup{(D.2.19y)--(D.2.19z)} extends through
source relative dimension nine. We keep that same
map and add the pullback of the correction along
$K_9\to Q_9$. This gives an actual ordinary map
$G_9:K_9\to T$ extending the allowed $G_8$ and
retaining every six-dimensional value in D.2.17.
The canonical summand has only been supplied
through dimension nine, and no extension over
the tenth or later cells follows from this proof.
